% Main LaTeX file for Causal Inference Notes
% Generated from aca-reader workflow

\documentclass[12pt]{amsbook}
\usepackage{amsmath, amssymb, amsthm, amsbsy, bm, mathtools}
\usepackage{xeCJK}
\usepackage{microtype}
\usepackage{hyperref}
\usepackage{cleveref}
\usepackage{natbib}
\usepackage{geometry}
\usepackage{graphicx}
\usepackage{float}
\usepackage{tikz}
\usetikzlibrary{arrows,positioning}
\geometry{margin=1in}

% Theorem styles - 标准定义环境
\theoremstyle{plain}
\newtheorem{Definition}{定义}[chapter]
\newtheorem{Theorem}[Definition]{定理}
\newtheorem{Lemma}[Definition]{引理}
\newtheorem{Corollary}[Definition]{推论}
\newtheorem{Proposition}[Definition]{命题}
\newtheorem{Example}{例}[chapter]
\newtheorem{Remark}{注}[chapter]
\newtheorem{Exercise}{练习}[chapter]
\newtheorem{Assumption}{假设}[chapter]
\newtheorem*{Proof}{证明}

% User annotation command
\newcommand{\userannotation}[2]{%
  \begin{Remark}
    \textbf{#1:} #2
  \end{Remark}
}

% Independence symbol (standard notation in causal inference)
\DeclareMathOperator{\Perp}{\mathrel{\perp\!\!\!\perp}}

% Title
\title{因果推断讲义}
\author{Peng Ding}
\date{\today}

\begin{document}

\maketitle
\tableofcontents

% Include chapters
% Chapter 1: Correlation, Association, and the Yule–Simpson Paradox
% A First Course in Causal Inference - Peng Ding

\chapter{相关性、关联与 Yule-Simpson 悖论}\label{ch:1}

\section{本章导论}

本章我们要回答一个看似简单却至关重要的问题：\textbf{相关性能否直接告诉我们因果关系？}

在日常生活中，我们经常听到这样的话：
\begin{itemize}
  \item "吸烟会导致肺癌"
  \item "教育水平越高，收入越高"
  \item "经常锻炼的人更长寿"
\end{itemize}

这些陈述都试图建立因果关系。但我们是如何得出这些结论的？仅仅因为我们观察到了相关性（correlation）就足够了么？

\section{为什么这个问题很重要}

想象一下以下场景：

\begin{Example}
一家公司想要知道某种新药是否有效。他们给一组患者服用新药，另一组服用安慰剂。结果显示服药组的康复率更高。这能否说明药物有效？
\end{Example}

直觉告诉我们，这似乎能说明问题。但仔细想想，还有很多其他因素可能影响结果：
\begin{itemize}
  \item 也许服药组的患者本身就更健康
  \item 也许医生在给药时，无意中对服药组患者更加关照
  \item 也许这只是随机波动
\end{itemize}

这就是因果推断（causal inference）的核心挑战：\textbf{我们观察到的关联（association），并不等于因果关系（causal relationship）}。

\section{相关性与因果性的区别}

统计学中，我们用不同的术语来区分这些概念：

\begin{itemize}
  \item \textbf{相关性（Correlation）}: 两个变量共同变化的趋势
  \item \textbf{关联（Association）}: 统计学意义上的依赖关系
  \item \textbf{因果性（Causation）}: 一个变量的变化直接导致另一个变量变化
\end{itemize}

关键洞察：\textit{相关性可以是因果的，也可以是非因果的。}

\section{二维列联表与关联度量}

将 $Z$ 视为治疗或暴露，$Y$ 视为结果，可以定义以下关联度量：

\[
\begin{tabular}{lcc}
\hline
        & $Y =1$ & $Y =0$\\\hline
$Z =1$ & $p_{11}$ & $p_{10}$\\
$Z =0$ & $p_{01}$ & $p_{00}$\\
\hline
\end{tabular}
\]

\textbf{风险差（Risk Difference）}：
\[
\RD = \mathbb{P}(Y=1 \mid Z=1) - \mathbb{P}(Y=1 \mid Z=0)
= \frac{p_{11}}{p_{11} + p_{10}} - \frac{p_{01}}{p_{01} + p_{00}}
\]

\textbf{风险比（Risk Ratio）}：
\[
\RR = \frac{\mathbb{P}(Y=1 \mid Z=1)}{\mathbb{P}(Y=1 \mid Z=0)}
= \frac{p_{11}}{p_{11} + p_{10}} \Big / \frac{p_{01}}{p_{01} + p_{00}}
\]

\textbf{优势比（Odds Ratio）}：
\[
\OR = \frac{\mathbb{P}(Y=1 \mid Z=1)/\mathbb{P}(Y=0 \mid Z=1)}{\mathbb{P}(Y=1 \mid Z=0)/\mathbb{P}(Y=0 \mid Z=0)}
= \frac{p_{11} p_{00}}{p_{10} p_{01}}
\]

\textbf{关于记号}：本书使用符号 $\Perp$ 表示随机变量的独立性或条件独立性。

\begin{Proposition}[二维列联表中的独立性]\label{prop:2x2-independence}
\begin{enumerate}
  \item 以下命题等价：$Z \Perp Y$、$\RD = 0$、$\RR = 1$、$\OR = 1$。
  \item 若所有 $p_{zy} > 0$，则 $\RD > 0$ 等价于 $\RR > 1$，也等价于 $\OR > 1$。
  \item 若 $\mathbb{P}(Y=1 \mid Z=1)$ 和 $\mathbb{P}(Y=1 \mid Z=0)$ 都很小，则 $\OR \approx \RR$。
\end{enumerate}
\end{Proposition}

\section{Yule-Simpson 悖论}

悖论的核心思想：在总体中观察到的关联方向，可能与分层后各组中的关联方向相反。

\begin{Example}[经典的 Yule-Simpson 悖论例子]
考虑一个大学录取数据：

\begin{center}
\begin{tabular}{c|c|c}
\hline
院系 & 男生录取率 & 女生录取率 \\
\hline
A & 62\% (341/550) & 82\% (137/167) \\
B & 63\% (28/44) & 68\% (24/35) \\
\hline
\textbf{总体} & \textbf{45\% (369/816)} & \textbf{30\% (161/525)} \\
\hline
\end{tabular}
\end{center}

有趣的现象：
\begin{itemize}
  \item 在 A 院系，男女生录取率分别是 62\% 和 82\%
  \item 在 B 院系，男女生录取率分别是 63\% 和 68\%
  \item 但总体来看，男生的录取率反而更高！
\end{itemize}

这就是 Yule-Simpson 悖论：\textbf{在每个分层中都有利的因素，在总体中可能不利}。
\end{Example}

这个例子揭示了一个深刻的问题：\textit{如果我们不控制混杂因素（confounding factors），就可能得出完全错误的因果结论}。

\section{偏相关系数与 Yule-Simpson 悖论}

考虑三维正态随机向量：
\[
\begin{pmatrix}
X \\ Y \\ Z
\end{pmatrix}
\sim N\left(
\begin{pmatrix}
0 \\ 0 \\ 0
\end{pmatrix},
\begin{pmatrix}
1 & \rho_{XY} & \rho_{XZ} \\
\rho_{XY} & 1 & \rho_{YZ} \\
\rho_{XZ} & \rho_{YZ} & 1
\end{pmatrix}
\right)
\]

$Y$ 和 $Z$ 的偏相关系数 $\rho_{YZ|X}$ 定义为在条件分布 $(Y, Z) \mid X$ 中的相关系数。

\begin{Proposition}[偏相关系数公式]\label{prop:partial-correlation}
\[
\rho_{YZ|X} = \frac{\rho_{YZ} - \rho_{YX}\rho_{ZX}}{\sqrt{1-\rho_{YX}^2}\sqrt{1-\rho_{ZX}^2}}
\]
\end{Proposition}

\textbf{重要含义}：即使 $\rho_{YZ} > 0$，偏相关系数 $\rho_{YZ|X}$ 仍可能为负！这正是正态随机向量情形下的 Yule-Simpson 悖论。

\section{LaLonde 数据与规格搜索问题}

LaLonde 数据集包含参与职业培训项目（National Supported Work Demonstration）的参与者的结果。在观察性研究中，协变量选择的敏感性是一个重要问题。

当我们运行 $2^{10} = 1024$ 种可能的协变量子集回归时，treatment 系数的显著性可能对模型规格高度敏感。这提醒我们：在观察性研究中，规格搜索可能导致误导性结论。

\section{种族歧视的简历实验}

Bertrand and Mullainathan (2004) 通过随机改变简历上的名字来研究种族歧视。这些名字被设计成听起来像非裔美国人或白人的名字。实验显示，有黑人名字的简历回调率显著较低。

当我们分别对男性和女性进行分析时，可能会发现有趣的亚组差异——这进一步说明"平均效应"可能掩盖了重要的异质性。

\section{本章小结与下节预告}

本章的核心要点：
\begin{enumerate}
  \item \textbf{相关性 $\neq$ 因果性}：观察到的关联可能是由于混杂因素造成的
  \item \textbf{关联度量}：$\RD$、$\RR$、$\OR$ 及它们与独立性的关系
  \item \textbf{Yule-Simpson 悖论}：总体关联可能与分层关联方向相反
  \item \textbf{偏相关}：即使边际相关为正，偏相关仍可能为负
  \item \textbf{因果推断的必要性}：我们需要系统性的方法来区分真实的因果效应
\end{enumerate}

下一章我们将学习因果推断的基石——\textbf{潜在结果框架（Potential Outcomes Framework）}，这将为我们的因果分析提供坚实的数学基础。

\section{习题}\label{sec:1-homework}

\begin{HomeworkProblem}[Independence in two-by-two tables]\label{exr:1-1}

Prove (1) and (2) in Proposition \ref{prop::2x2-independence}.


\end{HomeworkProblem}

\begin{HomeworkProblem}[More examples of the Yule--Simpson Paradox]\label{exr:1-2}


Give a numeric example of a two-by-two-by-two table in which the Yule--Simpson Paradox arises.

Find a real-life example in which the Yule--Simpson Paradox arises.



\end{HomeworkProblem}

\begin{HomeworkProblem}[Correlation and partial correlation]\label{exr:1-3}

Consider a three-dimensional Normal random vector:
\[
\begin{pmatrix}
X\\
Y\\
Z
\end{pmatrix}
\sim
\textsc{N}
\left(
\begin{pmatrix}
0\\
0\\
0
\end{pmatrix},
\begin{pmatrix}
1 & \rho_{XY} &\rho_{XZ} \\
  \rho_{XY}  & 1 & \rho_{YZ} \\
  \rho_{XZ} &  \rho_{YZ} & 1
\end{pmatrix}
\right) .
\]
The correlation coefficient between $Y$ and $Z$ is $\rho_{YZ}$.
There are many equivalent definitions of the partial correlation coefficient.
For a multivariate Normal vector, let $\rho_{YZ|X}$ denote the partial correlation coefficient between $Y$ and $Z$ given $X$, which is defined as their correlation coefficient in the conditional distribution $(Y, Z)\mid X$. Show that
\[
\rho_{YZ\mid X} = \frac{   \rho_{YZ} - \rho_{YX} \rho_{ZX}  }{  \sqrt{  1-\rho_{YX}^2 }  \sqrt{1- \rho_{ZX}^2  }  }  .
\]

Give a numerical example with $\rho_{YZ} > 0$ and $\rho_{YZ|X} < 0$.

Remark: This is the Yule--Simpson Paradox for a Normal random vector.
You can use the results in Chapter \ref{sec::multivariate-normal} to prove the formula for the partial correlation coefficient.



\end{HomeworkProblem}

\begin{HomeworkProblem}[Specification searches]\label{exr:1-4}

Section \ref{section:correlatioandregressionintro} re-analyzes the data used by \citet{hainmueller2012entropy}. We used the outcome named \ri{re78} and the treatment named \ri{treat} in the analysis. Moreover, the data also contain $10$ covariates and therefore $2^{10} = 1024$ possible subsets of covariates in the linear regression. Run $1024$ linear regressions with all possible subsets of covariates, and report the regression coefficients of the treatment. How many coefficients of the treatment are positively significant, how many are negatively significant, and how many are not significant? You can also report other interesting findings from these regressions.


\end{HomeworkProblem}

\begin{HomeworkProblem}[More on racial discrimination]\label{exr:1-5}

Section \ref{sec::twobytwotables-withresume} re-analyzes the data collected by \citet{bertrand2004emily}. Conduct analyses separately for males and females. What do you find from these subgroup analyses?


\end{HomeworkProblem}

\begin{HomeworkProblem}[Recommended reading]\label{exr:1-6}

\citet{bickel1975sex} is the original paper for the paradox reported in Section \ref{sec::berkeleydata}.
\citet{pearl2018book} recently revisited the study from the causal inference perspective.



\end{HomeworkProblem}

% Chapter 2: Potential Outcomes
% A First Course in Causal Inference - Peng Ding
% Rewritten in Stein motivation-first style with textbook formulas

\chapter{潜在结果框架}

\section{从相关到因果：为什么关联不够？}

在第一章中，我们见证了一个令人不安的现象：Yule-Simpson 悖论向我们展示了关联可以与因果相悖。两组数据分别显示的正向效应，在合并后可能变为负向——这警告我们，从关联推断因果是危险的。

但我们不禁要问：\textbf{关联毕竟是我们能观察到的，而因果是我们真正想知道的。那么，因果效应能否被清晰地定义，甚至被估计？}

这个问题引导我们走向一个优雅而强大的框架——\textbf{潜在结果框架（Potential Outcomes Framework）}，也被称为 \textbf{Neyman-Rubin 因果模型}。

这个框架的起源可以追溯到 Neyman 在 1923 年的开创性工作，以及 Rubin 在 1970 年代对其的进一步发展。它的核心思想其实出奇地简单：要谈论因果，我们必须想象（counterfactual）如果事情有所不同会发生什么。

\section{一个思想实验}

让我们从一个最简单的场景开始：假设我们想研究某种新药对患者康复的因果效应。

考虑第 $i$ 个患者。我们记：
\begin{itemize}
  \item $T_i = 1$：患者接受了治疗
  \item $T_i = 0$：患者未接受治疗（对照组）
\end{itemize}

现在，关键的问题来了：这个患者的结果会是什么？

在现实中，我们只能观察到一种情况——要么患者接受了治疗，要么没有。但为了谈论因果，我们必须同时考虑两种可能性：

\begin{itemize}
  \item $Y_i(1)$：\textbf{如果}患者接受治疗，会观察到的结果
  \item $Y_i(0)$：\textbf{如果}患者未接受治疗，会观察到的结果
\end{itemize}

这里，$Y_i(1)$ 和 $Y_i(0)$ 就是\textbf{潜在结果}——它们描述了在不同治疗状态下可能发生的结果，但我们永远无法同时观察到两者。

\begin{Remark}[注 2.1 — 什么是反事实]
\textbf{反事实}（counterfactual）指的是未被观测到的潜在结果。具体而言：
\begin{itemize}
  \item 对于治疗组个体 $i$（$T_i=1$），我们观察到 $Y_i(1)$，而 $Y_i(0)$ 是反事实
  \item 对于对照组个体 $i$（$T_i=0$），我们观察到 $Y_i(0)$，而 $Y_i(1)$ 是反事实
\end{itemize}
正因为潜在结果框架讨论的是``如果接受相反治疗会怎样''，所以它也被称为\textbf{反事实框架}（counterfactual framework）。
\end{Remark}

\userannotation{为什么叫``潜在''结果？}{因为它们描述的是``如果...会怎样''的情形，这种结果是``潜在的''而非``现实的''。这正是 counterfactual（反事实）思维的本质。}

\section{科学表（Science Table）}

\cite{rubin2005} 将潜在结果组成的 $n \times 2$ 矩阵称为\textbf{科学表（Science Table）}：

\begin{center}
\begin{tabular}{ccc}
\hline
$i$ & $Y_i(1)$ & $Y_i(0)$ \\
\hline
$1$ & $Y_1(1)$ & $Y_1(0)$ \\
$2$ & $Y_2(1)$ & $Y_2(0)$ \\
$\vdots$ & $\vdots$ & $\vdots$ \\
$n$ & $Y_n(1)$ & $Y_n(0)$ \\
\hline
\end{tabular}
\end{center}

这个表包含了所有我们需要知道的关于因果效应的信息——但问题是，我们永远只能观察到每一行的其中一个值。

\section{个体因果效应}

有了潜在结果，我们终于可以清晰地定义因果效应了。

对于个体 $i$，定义\textbf{个体因果效应（Individual Treatment Effect, ITE）}为：
\[
\tau_i = Y_i(1) - Y_i(0)
\]

这是治疗相对于对照的``净效果''。

但是，这里出现了一个根本性的困难。1923 年，Jerzy Neyman 在一篇关于田间试验的论文中首次明确指出了这个问题：

\begin{center}
\textit{对于每个个体，我们只能观察到一种潜在结果——实际发生的那种。另一种潜在结果是潜在的、永不可见的。}
\end{center}

这被称为\textbf{因果推断的根本问题（Fundamental Problem of Causal Inference）}。

\userannotation{为什么这个困难如此根本？}{因为科学上感兴趣的大多数因果问题，都涉及将现实与反事实进行比较。我们可以问``这种药对这个患者有效吗？''——但这个问题要求我们同时知道患者吃药和不吃药的结果，而这是不可能的。}

\section{平均因果效应}

既然个体因果效应不可测量，我们退而求其次，关注\textbf{平均因果效应（Average Causal Effect, ACE）}：
\[
\tau = \mathbb{E}[Y_i(1) - Y_i(0)] = \bar{Y}(1) - \bar{Y}(0)
\tag{2.1}\label{eq:ace}
\]

用更精确的有限总体记号：\footnote{推导见附录 \cref{sec:appendix-ace-finite-sample}}
\[
\tau = \frac{1}{n} \sum_{i=1}^n \{ Y_i(1) - Y_i(0) \} = \frac{1}{n} \sum_{i=1}^n Y_i(1) - \frac{1}{n} \sum_{i=1}^n Y_i(0)
\tag{2.2}\label{eq:ace-finite}
\]

\textbf{一个关键恒等式}：平均因果效应等于个体因果效应的均值：
\[
\tau = \bar{Y}(1) - \bar{Y}(0) = \frac{1}{n} \sum_{i=1}^n \tau_i
\tag{2.3}\label{eq:ace-ite-mean}
\]

\section{亚组与亚组因果效应}

\textbf{亚组（subgroup）}是由某个协变量（通常是二元变量 $X_i$）分割出来的子群体。例如：
\begin{itemize}
  \item 按性别分割：男性亚组（$X_i=1$）和女性亚组（$X_i=0$）
  \item 按年龄分割：老年人/年轻人
  \item 按收入分割：高收入/低收入
\end{itemize}

\textbf{亚组因果效应（Subgroup Causal Effect）}衡量的是在该亚组内的平均因果效应：\footnote{推导见附录 \cref{sec:appendix-subgroup-ace}}
\[
\tau_x = \frac{\sum_{i=1}^n \mathbb{I}(X_i = x) \{ Y_i(1) - Y_i(0) \}}{\sum_{i=1}^n \mathbb{I}(X_i = x)}, \quad (x = 0, 1)
\tag{2.4}\label{eq:subgroup-ace}
\]

\section{亚组因果效应与 Yule-Simpson 悖论的不存在}

一个重要的恒等式揭示了亚组因果效应与总体因果效应的关系：\footnote{推导见附录 \cref{sec:appendix-subgroup-decomposition}}
\[
\tau = \pi_1 \tau_1 + \pi_0 \tau_0
\tag{2.5}\label{eq:ace-subgroup-decomposition}
\]
其中 $\pi_x = \frac{1}{n} \sum_{i=1}^n \mathbb{I}(X_i = x)$ 是 $X_i = x$ 的单元比例。

\textbf{关键推论}：如果 $\tau_1 > 0$ 且 $\tau_0 > 0$，则必有 $\tau > 0$。因此，\textbf{Yule-Simpson 悖论在因果效应层面不可能发生！}

这澄清了第一章的困惑：在关联层面出现的悖论，在因果层面并不存在——因为因果效应是可分解的，而关联并非如此。

\section{治疗分配机制与基本恒等式}

设 $T_i$ 是单元 $i$ 的二元治疗指示变量，观测结果 $Y_i$ 是潜在结果和治疗分配的函数：

\begin{align}
Y_i &= \begin{cases}
Y_i(1), & \text{if } T_i = 1 \\
Y_i(0), & \text{if } T_i = 0
\end{cases} \label{eq:observed-y} \\
&= T_i Y_i(1) + (1 - T_i) Y_i(0) \label{eq:fundamental-bridge} \\
&= Y_i(0) + T_i \{ Y_i(1) - Y_i(0) \} \label{eq:treatment-effect} \\
&= Y_i(0) + T_i \tau_i \label{eq:treatment-effect-2}
\end{align}

\cite{pearl2010} 将 \cref{eq:fundamental-bridge} 视为潜在结果与观测结果之间的\textbf{基本桥梁}。这些方程强调了个体因果效应 $\tau_i = Y_i(1) - Y_i(0)$ 可能在不同单元间存在异质性。

\section{SUTVA：使框架有效的必要假设}

使用潜在结果框架需要一些技术性假设。其中最重要的是\textbf{稳定单位治疗值假设（Stable Unit Treatment Value Assumption, SUTVA）}。

\begin{Assumption}[无干扰]\label{assump:sutva1}
单元 $i$ 的潜在结果不依赖于其他单元的治疗分配。这有时被称为\textbf{无干扰}假设。
\end{Assumption}

\begin{Assumption}[一致性]\label{assump:sutva2}
治疗水平没有其他版本。等价地，我们要求治疗水平是良好定义的，至少对关注的结局是如此。这有时被称为\textbf{一致性}假设。
\end{Assumption}

\begin{Assumption}[SUTVA]\label{assump:sutva-complete}
假设 \ref{assump:sutva1} 和 \ref{assump:sutva2} 同时成立。
\end{Assumption}

\textbf{为什么 SUTVA 如此重要？} 如果单元之间存在干扰，或者治疗定义不明确，那么潜在结果 $Y_i(1)$ 和 $Y_i(0)$ 就不是well-defined的——它们的值将取决于其他人的行为或治疗的具体形式，整个框架就会崩塌。

\textbf{违反 SUTVA 的例子}：

\textbf{无干扰假设（no-interference assumption）违反的情形}：

\begin{enumerate}
  \item \textbf{传染病与疫苗接种（infectious diseases and vaccination）}：当疾病具有传染性时，一个人是否接种疫苗不仅影响他自己的患病风险，还影响他周围人被感染的概率。例如，如果我的同事接种了流感疫苗，我患流感的概率也会降低——即使我自己没有接种。在这种情况下，单元 $i$ 的潜在结果 $Y_i(1)$ 和 $Y_i(0)$ 都依赖于其他人的治疗分配。

  \item \textbf{网络效应实验（network experiments）}：在社交网络实验中，如果一个人收到 Facebook 广告（treatment），他的朋友看到这条广告并受到影响后，可能也会购买该产品——即使朋友本人没有直接收到广告。这就是所谓的\textbf{溢出效应（spillover effect）}。

  \item \textbf{随机区组实验中的相邻地块（neighboring plots in randomized block experiments）}：农业实验中，相邻地块的作物可能相互影响——如果相邻地块使用了肥料，可能会通过土壤渗透影响本地块的产量。
\end{enumerate}

\textbf{一致性假设（consistency assumption）违反的情形}：

\begin{enumerate}
  \item \textbf{复合治疗（composite treatments）}：当治疗本身有复杂组成部分时，相同的``治疗''标签可能对应不同的实际干预。例如：
  \begin{itemize}
    \item 研究吸烟（smoking）对肺癌的影响时，``吸烟''这个治疗有多种形式——香烟类型、吸烟频率、吸入深度等，这些都可能导致不同的潜在结果。
    \item 研究大学教育（college education）对收入的影响时，``上大学''这个治疗包括大学类型（985/211 vs 普通院校）、专业、文凭等级等，这些差异会造成潜在结果的不一致。
  \end{itemize}

  \item \textbf{剂量效应（dosage effects）}：在医学研究中，相同的药物名称可能对应不同剂量，而剂量直接影响疗效。例如``服用阿司匹林''作为治疗，但不同剂量的阿司匹林会产生完全不同的潜在结果。

  \item \textbf{实施变异（implementation variation）}：即使治疗名义上相同，不同实施者可能引入差异。例如，心理治疗研究中，不同治疗师的风格和技能可能导致相同的``心理治疗''产生不同效果。
\end{enumerate}

当 SUTVA 不成立时，潜在结果框架需要修正。\textbf{现代因果推断文献}（如 \cite{halloran1996}）正在积极研究如何处理存在干扰或治疗版本多样的情形。

\userannotation{SUTVA 的历史重要性}{Rubin (1980) 将 SUTVA 确立为潜在结果框架的基础假设，没有 SUTVA，因果推断的数学化是不可能的。}

\section{反事实思维的逻辑}

也许你会好奇：为什么不能简单地问``实际发生的结果是什么''？

答案是：\textbf{因果效应本质上是一个反事实（counterfactual）概念}。当我们问``这种药是否有效''时，我们实际上是在问：

\begin{center}
\textit{``如果这个患者吃了药，结果会怎样？但现实是，他可能没有吃药...''}
\end{center}

这种思维方式可能让人觉得抽象甚至哲学化，但它是我们谈论因果的唯一严谨方式。潜在结果框架的价值在于，它将这种哲学洞察转化为清晰的数学语言。

\section{实验单元定义的微妙性}

我讨论一个与实验单元定义相关的微妙之处。简单来说，\textbf{实验单元可能不同于物理单元}。

例如，如果我在服药前没有服用阿司匹林且头痛没有消失，但我现在服用阿司匹林后头痛消失了，你可能认为我们可以同时观察到我的两种潜在结果。但这是对实验单元定义的误解！在不同时刻，我是同一个物理人，但成为了两个不同的实验单元。

因此，我们有四个潜在结果，其中两个被观察到，两个缺失。这说明，即使观察到的数据显示治疗``有效''，我们也无法确定地说阿司匹林就是原因——因为在``治疗后''的状态下，可能存在其他因素（如时间效应）导致头痛消失。

\section{非线性因果参数}

平均因果效应是一个\textbf{线性}因果参数，因为潜在结果的平均差等于个体差别的平均。其他参数可能不是线性的。

例如，中位数治疗效应：
\[
\delta_1 = \text{median}\{ Y_i(1) \}_{i=1}^n - \text{median}\{ Y_i(0) \}_{i=1}^n
\tag{2.6}\label{eq:median-treatment-effect-1}
\]
通常不同于个体治疗效应的中位数：
\[
\delta_2 = \text{median}\{ Y_i(1) - Y_i(0) \}_{i=1}^n
\tag{2.7}\label{eq:median-treatment-effect-2}
\]

这两个量在某些情况下可能相差很大。\footnote{数值例子见附录 \cref{sec:appendix-nonlinear-causal}}

\section{本章小结}

本章建立因果推断的数学语言：

\begin{enumerate}
  \item \textbf{潜在结果}：为每个个体定义 $Y_i(1)$ 和 $Y_i(0)$
  \item \textbf{科学表}：$n \times 2$ 矩阵 $\{Y_i(1), Y_i(0)\}_{i=1}^n$
  \item \textbf{个体因果效应}：$\tau_i = Y_i(1) - Y_i(0)$
  \item \textbf{平均因果效应}：$\tau = \bar{Y}(1) - \bar{Y}(0) = n^{-1}\sum_{i=1}^n \tau_i$
  \item \textbf{根本问题}：我们只能观察到一种潜在结果
  \item \textbf{基本桥梁方程}：$Y_i = T_i Y_i(1) + (1-T_i) Y_i(0)$
  \item \textbf{SUTVA}：一致性 + 无干扰
  \item \textbf{Yule-Simpson 悖论}在因果框架下不可能发生
\end{enumerate}

下一章我们将看到，\textbf{随机实验}提供了一种绕过根本问题的优雅方式——通过随机分配治疗，我们可以建立组间可比性，从而识别因果效应。

\section{习题}\label{sec:chapter2-exercises}

\begin{Exercise}{\ref{ex:2.1} 完美的医生}\label{ex:2.1}
假设潜在结果由以下模型生成：
\[
Y_i(0) \sim N(0, 1), \quad \tau_i = -0.5 + Y_i(0), \quad Y_i(1) = Y_i(0) + \tau_i
\]
治疗分配为 $T_i = \mathbb{I}(\tau_i \geq 0)$，观测结果为 $Y_i = T_i Y_i(1) + (1-T_i) Y_i(0)$。

计算条件均值差 $\mathbb{E}(Y_i \mid T_i=1) - \mathbb{E}(Y_i \mid T_i=0)$。

\textbf{提示}：需要使用截断正态分布的均值公式。当 $X \sim N(\mu, \sigma^2)$ 时，
\[
\mathbb{E}(X \mid a < X < b) = \mu - \sigma \frac{\phi\left(\frac{b-\mu}{\sigma}\right) - \phi\left(\frac{a-\mu}{\sigma}\right)}{\Phi\left(\frac{b-\mu}{\sigma}\right) - \Phi\left(\frac{a-\mu}{\sigma}\right)}
\]
其中 $\phi$ 和 $\Phi$ 分别是标准正态分布的密度函数和分布函数。
\end{Exercise}

\begin{Exercise}{\ref{ex:2.2} 非线性因果参数}\label{ex:2.2}
证明平均因果效应是线性因果参数的例子，即 $\bar{Y}(1) - \bar{Y}(0) = n^{-1}\sum_{i=1}^n \{Y_i(1) - Y_i(0)\}$。

其他参数可能不是线性的。例如，定义中位数治疗效应为：
\[
\delta_1 = \text{median}\{Y_i(1)\} - \text{median}\{Y_i(0)\}
\]
这通常不同于个体治疗效应的中位数：
\[
\delta_2 = \text{median}\{Y_i(1) - Y_i(0)\}
\]

\begin{enumerate}
  \item 给出数值例子分别满足 $\delta_1 = \delta_2$、$\delta_1 > \delta_2$ 和 $\delta_1 < \delta_2$。
  \item 哪个参数更有意义？为什么？请用例子来证明你的结论。
\end{enumerate}
\end{Exercise}

\begin{Exercise}{\ref{ex:2.3} 平均效应与个体效应}\label{ex:2.3}
给出一个数值例子，其中 $\tau = n^{-1}\sum_{i=1}^n \{Y_i(1) - Y_i(0)\} > 0$，但有不到一半的单元满足 $Y_i(1) > Y_i(0)$。即平均因果效应为正，但治疗对不到一半的单元有益。
\end{Exercise}

\begin{Exercise}{\ref{ex:2.4} 推荐阅读}\label{ex:2.4}
\begingroup
\parskip=0pt
\textbf{推荐阅读}：\cite{holland1986} 是统计因果推断领域的经典综述文章，它推广了``Rubin 因果模型''这一名称。在 UC Berkeley，我们称之为``Neyman 模型''——原因显而易见。
\endgroup
\end{Exercise}


% Chapter 3: Randomized Experiments
% A First Course in Causal Inference - Peng Ding
% Rewritten in Stein motivation-first style with textbook formulas

\chapter{完全随机实验与费舍尔随机化检验}

\section{从潜在结果到随机实验}

在上一章中，我们建立了潜在结果框架，学会了如何清晰地定义因果效应。平均因果效应定义为：
\[
\tau = \mathbb{E}[Y_i(1) - Y_i(0)] = \bar{Y}(1) - \bar{Y}(0)
\]

但我们也清楚地看到了因果推断的根本问题：对于每个个体，我们只能观察到一种潜在结果，$Y_i(1)$ 和 $Y_i(0)$ 无法同时被观测。

那么，一个自然的问题是：我们能否设计某种数据收集过程，使得即使无法观察反事实，我们仍然可以估计因果效应？

这个问题的答案是肯定的，而解决方案就是\textbf{随机实验（Randomized Experiment）}——统计学的黄金标准。

\section{Fisher 的洞见：随机化的两个原则}

Ronald Fisher 在 1935 年的开创性著作《Design of Experiments》中指出了随机化的两个优势：

\begin{enumerate}
  \item \textbf{可比性}：它创造了在平均意义上可比的 treatment 和 control 组。
  \item \textbf{推理基础}：它作为统计推断的"合理基础"。
\end{enumerate}

第一点直观易懂——随机治疗分配不会偏向 treatment 或 control。第二点更微妙：Fisher 的意思是，随机化为统计检验提供了正当理由，这就是现在所谓的\textbf{费舍尔随机化检验（Fisher Randomization Test, FRT）}。

\section{完全随机实验的定义}

考虑一个包含 $n$ 个单元的实验，其中 $n_1$ 个接受治疗，$n_0$ 个接受对照（$n = n_1 + n_0$）。

\begin{Definition}[CRE]\label{def:cre}
固定 $n_1$ 和 $n_0$（其中 $n = n_1 + n_0$）。完全随机实验的治疗 assignment 机制为：
\[
\mathbb{P}(\bm{Z} = \bm{z}) = 1 \Big / \binom{n}{n_1},
\]
其中 $\bm{z} = (z_1,\ldots, z_n)$ 满足 $\sumn z_i = n_1$ 和 $\sumn (1-z_i) = n_0$。
\end{Definition}

在这个定义中，我们假设潜在结果向量 $\bm{Y}(1) = (Y_1(1), \ldots, Y_n(1))$ 和 $\bm{Y}(0) = (Y_1(0), \ldots, Y_n(0))$ 都是固定的。即使我们将它们视为随机的，我们可以在其上条件化，治疗分配机制变为：
\[
\mathbb{P}\{\bm{Z} = \bm{z} \mid \bm{Y}(1), \bm{Y}(0)\} = 1 \Big / \binom{n}{n_1}
\]
因为在 CRE 中 $\bm{Z} \ind \{\bm{Y}(1), \bm{Y}(0)\}$。

\textbf{关键含义}：在 CRE 中，$\bm{Z}$ 是从 $n_1$ 个 1 和 $n_0$ 个 0 的随机排列中产生的。

\section{Sharp Null 假设与 FRT 逻辑}

\citet{Fisher:1935} 感兴趣检验的零假设是：
\[
\HF: Y_i(1) = Y_i(0) \text{ 对所有单元 } i=1,\ldots, n.
\]
\citet{Rubin:1980} 称其为\textbf{sharp null 假设}，因为它可以基于观测数据确定所有潜在结果：$\bm{Y}(1) = \bm{Y}(0) = \bm{Y} = (Y_1, \ldots, Y_n)$，即观测结果向量。它也被称为\textbf{强零假设}。

概念上，在 $\HF$ 下，FRT 适用于任何检验统计量：
\begin{equation}\label{eq:teststatistic}
T = T(\bm{Z}, \bm{Y}) ,
\end{equation}
它是观测数据的函数。观测结果向量 $\bm{Y}$ 在 $\HF$ 下是固定的，所以 $T$ 中唯一的随机成分是治疗向量 $\bm{Z}$。实验者决定 $\bm{Z}$ 的分布，进而决定 $T$ 的分布——这就是计算 p 值的基础。

在 CRE 中，$\bm{Z}$ 在集合 $\{\bm{z}^1, \ldots, \bm{z}^M\}$ 上是均匀分布的，其中 $M = \binom{n}{n_1}$，$\bm{z}^m$ 是所有可能的包含 $n_1$ 个 1 和 $n_0$ 个 0 的向量。

\section{随机化分布与精确 p 值}

如果较大的 $T$ 值更极端，我们可以使用以下尾概率来衡量检验统计量相对于其随机化分布的极端性：
\begin{equation}\label{eq:exact-pvalue}
p_\frt = M^{-1} \sum_{m=1}^M \mathbb{I}\{T(\bm{z}^m,\bm{Y}) \geq T(\bm{Z},\bm{Y}) \}.
\end{equation}

这就是 Fisher 的 p 值。重要的是，$p_\frt$ 在 \cref{eq:exact-pvalue} 中对任何检验统计量和任何结果生成过程都有效。它在 $\HF$ 下是有限样本精确的（finite-sample exact）：
\begin{equation}\label{eq:p-value-valid}
\mathbb{P}(p_\frt \leq u) \leq u \quad \text{for all} \quad 0 \leq u \leq 1.
\end{equation}

\userannotation{为什么叫"有限样本精确"？}{这意味着 p 值的分布（虽然在离散情况下是阶梯函数）满足上式，因此 p 值在 $[0,1]$ 上均匀分布（对于连续情况）。}

\section{Monte Carlo 近似}

在实践中，$M$ 通常太大（例如当 $n=100, n_1=50$ 时，$M > 10^{29}$），枚举所有可能的治疗向量在计算上是不可行的。我们经常用 Monte Carlo 来近似 $p_\frt$：
\begin{equation}\label{eq:mc-pvalue}
\hat p_\frt  =  R^{-1} \sum_{r=1}^R   \mathbb{I}\{T(\bm{z}^r,\bm{Y})   \geq T(\bm{Z},\bm{Y})   \},
\end{equation}
其中 $\bm{z}^r$ 是 $\bm{Z}$ 的 $R$ 次随机置换。由于计算 $p_\frt$ 涉及对 $\bm{Z}$ 的置换，FRT 有时在 CRE 语境下被称为\textbf{置换检验（permutation test）}。

\textbf{修正的有限样本有效近似}：
\[
\tilde p_\frt  =  \frac{  1 + \sum_{r=1}^R   \mathbb{I}\{T(\bm{z}^r,\bm{Y})   \geq T(\bm{Z},\bm{Y})   \}  }{1+R}
\]
这个修正在任意有限的 $R$ 下都是有限样本有效的 p 值。

\section{检验统计量的规范选择}

FRT 的一个特点是它对任何检验统计量都生成有限样本精确 p 值。但我们应该选择能够给出 $\HF$ 可能违背信息的检验统计量。

\subsection{差值均值统计量}

差值均值统计量是：
\[
\hat{\tau}   = \hat{\bar{Y}}(1) - \hat{\bar{Y}}(0)
\]
其中
\[
\hat{\bar{Y}}(1)  =  n_1^{-1} \sum_{Z_i=1} Y_i  = n_1^{-1} \sumn Z_iY_i
\]
是 treatment 组的样本均值，
\[
\hat{\bar{Y}}(0) = n_0^{-1}  \sum_{Z_i=0} Y_i   = n_0^{-1}  \sumn (1-Z_i) Y_i
\]
是 control 组的样本均值。

在 $\HF$ 下，它有均值 $0$ 和方差：
\[
\text{var}(\hat{\tau})  = \frac{n}{n_1 n_0} s^2,
\]
其中
\[
s^2 = (n-1)^{-1} \sumn (Y_i - \bar{Y})^2, \quad \bar{Y} = n^{-1} \sumn Y_i.
\]

\textbf{渐近分布}：由于有限总体中心极限定理，随机化分布近似正态：
\begin{equation}\label{eq:frt-tau-hat}
\frac{ \hat{\tau} }{ \sqrt{\frac{n}{n_1 n_0} s^2} } \rightarrow \mathcal{N}(0,1)
\end{equation}
in distribution.

\textbf{与经典 t 检验的联系}：在 $\HF$ 下，近似 p 值接近经典两样本 t 检验（假设等方差）的 p 值。

基于代数运算（见习题 \ref{ex:3.8}），我们有如下展开式：
\begin{equation}\label{eq:frt-algebra}
(n-1)s^2 = \sum_{Z_i=1} \{Y_i - \hat{\bar{Y}}(1)\}^2 + \sum_{Z_i=0} \{Y_i - \hat{\bar{Y}}(0)\}^2 + \frac{n_1 n_0}{n} \hat{\tau}^2 .
\end{equation}

\userannotation{为什么 FRT 能证明 t 检验的合理性？}{传统 t 检验需要正态性假设。但 FRT 框架表明，只要随机化分配正确执行，t 检验程序在有限总体 + 随机化情境下可以使用——渐近理论只依赖于有限总体 CLT，而不依赖于数据的原始分布。}

\subsection{学生化统计量}

差值均值隐含地假设两组方差相等。如果治疗可能影响结果的变异性（异方差），我们需要使用学生化统计量：
\[
t = \frac{ \hat{\bar{Y}}(1) - \hat{\bar{Y}}(0) }{ \sqrt{ \frac{ \hat{S}^2(1) }{n_1} + \frac{\hat{S}^2(0)}{n_0} } },
\]
其中
\[
\hat{S}^2(1) = (n_1 - 1)^{-1} \sum_{Z_i =1} \{ Y_i -\hat{\bar{Y}}(1) \}^2,\quad
\hat{S}^2(0) = (n_0 - 1)^{-1} \sum_{Z_i =0} \{ Y_i -\hat{\bar{Y}}(0) \}^2
\]
分别是 treatment 组和 control 组的样本方差。

在 $\HF$ 下，有限总体 CLT 再次意味着 $t$ 渐近正态：$t \rightarrow \mathcal{N}(0,1)$。

这个统计量与 \texttt{t.test(var.equal = FALSE)} 的 p 值接近。

\userannotation{学生化为什么更稳健？}{当 treatment 组和 control 组的方差差异较大时，差值均值的检验可能失效。学生化通过分别估计两组的方差，能够更好地处理异方差情况。}

\subsection{Wilcoxon 秩和统计量}

差值均值和学生化统计量都关注均值差异，且对异常值敏感。Wilcoxon 秩和统计量基于合并观测结果的秩：
\[
W = \sumn Z_i R_i,
\]
其中 $R_i$ 是 $Y_i$ 在合并样本中的秩。

在 $\HF$ 下，$W$ 有均值 $\frac{n_1(n+1)}{2}$ 和方差 $\frac{n_1 n_0 (n+1)}{12}$，且渐近正态。

这个统计量的优点是：
\begin{itemize}
  \item 对异常值稳健
  \item 只依赖于数据的秩次信息
  \item 在许多备择假设下都有较好的功效
\end{itemize}

\subsection{Kolmogorov-Smirnov 统计量}

上述统计量主要关注均值差异。KS 统计量衡量两个分布之间的最大距离：
\[
D = \max_y |\hat{F}_1(y) - \hat{F}_0(y)|,
\]
其中 $\hat{F}_1(y) = n_1^{-1} \sumn Z_i I(Y_i \leq y)$。

这个统计量能够检测分布形状的差异，而不仅仅是均值差异。

\section{FRT 与传统 t 检验的关系}

一个重要的理论发现是：

\begin{center}
\textbf{FRT 为传统 t 检验提供了更坚实的基础。}
\end{center}

传统 t 检验假设"数据来自正态分布的无限总体"。但 FRT 框架表明，只要随机化分配正确执行，t 检验程序可以在"有限总体 + 随机化"的情境下使用。

这意味着，在随机实验中，我们可以不那么担心正态性假设——CLT 会为我们处理这个问题。

\section{LaLonde 实验数据案例}\label{sec:lalonde-data}

1986 年，Robert LaLonde 发表了一项影响深远的研究，探讨职业培训项目（National Supported Work Demonstration）对参与者收入的影响。

使用 R 进行 FRT 分析：

\begin{verbatim}
library(Matching)
data(lalonde)
z <- lalonde$treat
y <- lalonde$re78
\end{verbatim}

观测到的检验统计量值：

\begin{center}
\begin{tabular}{c|c|c}
\hline
统计量 & 观测值 & 描述 \\
\hline
$\hat{\tau}$ & 2.84 & 差值均值（t 统计量）\\
$t$ & 2.67 & 学生化统计量 \\
$W$ & 27402.5 & Wilcoxon 秩和 \\
$D$ & 0.132 & KS 统计量 \\
\hline
\end{tabular}
\end{center}

通过随机置换治疗向量，我们可以得到检验统计量的随机化分布的 Monte Carlo 近似。

Monte Carlo 置换检验的结果（$R = 10^4$）：

\begin{verbatim}
> exact.pv = c(mean(Tauhat >= tauhat),
+               mean(Student >= student),
+               mean(Wilcox >= W),
+               mean(Ks >= D))
> round(exact.pv, 3)
[1] 0.002 0.002 0.006 0.040
\end{verbatim}

所有单侧 p 值都小于 0.05——这提供了治疗效应显著的强有力证据。

\section{历史侧边栏：随机实验的起源}\label{sec:tea-history}

FRT 的逻辑可以追溯到早期的历史实例。

\subsection{James Lind 的坏血病实验}

James Lind (1716-1794) 是苏格兰医生和皇家海军卫生学的先驱。他进行了最早有清晰文献记录的随机实验之一，研究柑橘类水果对坏血病的影响。在 12 名坏血病患者被分配到 6 组后，接受柑橘类水果的组在六天后康复，而其他组没有。

如果将治疗简化为二元指示变量，观测到的极端 $2 \times 2$ 表格在 $\HF$ 下出现的概率为：
\[
p_\frt = \frac{1}{\binom{12}{2}} = \frac{1}{66} = 0.015.
\]

\subsection{Fisher 的女士品茶实验}

\citet{Fisher:1935} 描述了一个著名实验，一位女士声称能区分先加茶还是先加牛奶制成的奶茶。他制作了 8 杯茶（4 杯先加茶，4 杯先加牛奶），随机顺序呈现给女士。在她的全部猜对（$X = 4$）的情况下，p 值为：
\[
p_\frt = \frac{1}{70} = 0.014.
\]

这些例子强调了 Fisher 的两个原则：
\begin{enumerate}
  \item \textbf{随机化}：实验的随机分配保证了 p 值的有效性
  \item \textbf{重复}：必须有足够的单元以确保 FRT 有功效
\end{enumerate}

\section{Sharp Null 假设的置信区间}

FRT 的逻辑不仅限于 $\HF$。我们还可以检验更一般的 sharp null 假设：
\[
H_0(\bm{\tau}): Y_i(1) - Y_i(0) = \tau_i \quad \text{对所有 } i.
\]

通过检验与置信集的对偶性，我们可以构建 $\tau$ 的 $(1-\alpha)$ 置信集。

\section{本章小结}

本章建立了随机实验的理论与实践：

\begin{enumerate}
  \item \textbf{CRE 定义}：$\mathbb{P}(\bm{Z} = \bm{z}) = 1 / \binom{n}{n_1}$
  \item \textbf{随机化的两个原则}：可比性 + 推理基础
  \item \textbf{Sharp null}：$\HF: Y_i(1) = Y_i(0)$ 对所有 $i$
  \item \textbf{精确 p 值}：$p_\frt = M^{-1} \sum_{m=1}^M \mathbb{I}\{T(\bm{z}^m, \bm{Y}) \geq T(\bm{Z}, \bm{Y})\}$
  \item \textbf{多种统计量}：差值均值 $\hat{\tau}$、学生化 $t$、Wilcoxon $W$、KS $D$
  \item \textbf{理论联系}：FRT 为传统 t 检验提供基础
  \item \textbf{历史起源}：Lind 的坏血病实验、Fisher 的女士品茶
\end{enumerate}

下一章我们将学习 Neymanian 推断——另一种基于渐近理论的方法，它可以更方便地构建置信区间。

\section{习题}

\begin{Exercise}{3.1 — $p_\frt$ 的精确性}
\label{ex:3.1}
证明在 sharp null 下，$p_\frt$ 满足：
\[
\mathbb{P}(p_\frt \leq u) \leq u \quad \text{for all} \quad 0 \leq u \leq 1.
\]
\end{Exercise}

\begin{Exercise}{3.2 — Monte Carlo 误差}
\label{ex:3.2}
设 $p_\frt$ 是固定值，$\hat p_\frt$ 是其 Monte Carlo 估计（见 \cref{eq:mc-pvalue}）。证明：
\[
\mathbb{E}_{\text{mc}} (\hat p_\frt) = p_\frt, \quad
\text{var}_{\text{mc}} (\hat p_\frt) \leq \frac{1}{4R},
\]
其中下标 "mc" 表示由于 Monte Carlo 的随机性。
\end{Exercise}

\begin{Exercise}{3.3 — 有限样本有效的 Monte Carlo 近似}
\label{ex:3.3}
证明修正的 Monte Carlo 近似
\[
\tilde p_\frt  =  \frac{ 1 + \sum_{r=1}^R \mathbb{I}\{T(\bm{z}^r,\bm{Y}) \geq T(\bm{Z},\bm{Y}) \} }{1+R}
\]
在任意有限的 $R$ 下都是有限样本有效的 p 值。

\textbf{提示}：可以利用以下两个基本概率结果：
\begin{enumerate}
  \item 对于两个二项分布 $X_1 \sim \text{Binomial}(R, p_1)$ 和 $X_2 \sim \text{Binomial}(R, p_2)$，若 $p_1 \geq p_2$，则 $\mathbb{P}(X_1 \leq x) \leq \mathbb{P}(X_2 \leq x)$。
  \item 若 $p \sim \text{Uniform}(0,1)$ 且 $X \mid p \sim \text{Binomial}(R, p)$，则 $X$ 在 $\{0, 1,\ldots, R\}$ 上均匀分布。
\end{enumerate}
\end{Exercise}

\begin{Exercise}{3.4 — Fisher 精确检验}
\label{ex:3.4}
考虑二元结局的 CRE，数据汇总为 $2 \times 2$ 列联表：

\begin{center}
\begin{tabular}{lccc}
\hline
        & $Y =1$ & $Y =0$& 合计\\\hline
$Z =1$ & $n_{11}$ & $n_{10}$ & $n_1$\\
$Z =0$ & $n_{01}$ & $n_{00}$ & $n_0$\\
\hline
\end{tabular}
\end{center}

在 $\HF$ 下，证明任何检验统计量 $T(n_{11}, n_{10}, n_{01}, n_{00})$ 是 $n_{11}$ 的函数，且 $n_{11}$ 的精确分布是超几何分布。请说明超几何分布的参数。
\end{Exercise}

\begin{Exercise}{3.5 — 女士品茶详细计算}
\label{ex:3.5}
回忆女士品茶实验（见 \cref{sec:tea-history}）。计算 $\mathbb{P}(X = k)$，其中 $X$ 是女士正确识别 $k$ 杯先加牛奶的茶的个数，$k = 0, 1, 2, 3, 4$。
\end{Exercise}

\begin{Exercise}{3.6 — 协变量调整的 FRT}
\label{ex:3.6}
使用 LaLonde 数据集（见 \cref{sec:lalonde-data}），添加协变量进行 FRT 分析。假设所有潜在结果和协变量都是固定数值。

有两种策略：
\begin{enumerate}
  \item 运行结局对协变量的线性回归，得到残差（即"伪结局"）。基于残差定义四种检验统计量，进行 FRT 并报告 p 值。
  \item 定义检验统计量为结局对治疗和协变量线性回归中治疗系数的 FRT。
\end{enumerate}

解释为什么上述两种策略的五个 p 值都是有限样本精确的。
\end{Exercise}

\begin{Exercise}{3.7 — FRT 与广义线性模型}
\label{ex:3.7}
使用与 \cref{ex:3.6} 相同的数据集，但将结局改为 \texttt{re78 > 0} 的二元指示变量。运行结局对治疗和协变量的 logistic 回归。治疗系数是否显著？FRT p 值是多少？
\end{Exercise}

\begin{Exercise}{3.8 — 代数细节}
\label{ex:3.8}
验证 \cref{eq:frt-algebra}。
\end{Exercise}

\begin{Exercise}{3.9 — 推荐阅读}
\label{ex:3.9}
\citet{bind2020possible} 是一篇倡导在复杂实验中同时报告 p 值及其相应随机化分布的论文。
\end{Exercise}

% === 用户问答记录 ===
% (后续通过 interactive-qa 更新)

% Chapter 4: Neymanian Repeated Sampling Inference
% A First Course in Causal Inference - Peng Ding
% Rewritten in Stein motivation-first style with textbook theorems

\chapter{Neymanian 重复抽样推断}\label{ch:4}

\section{两种推断哲学的对话}

在上一章中，我们学习了费舍尔随机化检验（FRT）。Fisher 的方法优雅而精确——它基于随机化的精确性质，在 sharp null 假设下给出精确的 p 值。

但 FRT 也有其局限：它只能检验 sharp null 假设，对于构建置信区间需要额外的推理。而且当样本量很大时，枚举所有可能的置换变得计算上不可行。

Jerzy Neyman 在 1923 年的工作中提出了一种互补的方法。这种方法基于\textbf{渐近正态分布}，可以更方便地构建置信区间和进行假设检验。Neyman 的方法不依赖于任何特定的 null 假设，而是直接对因果效应进行统计推断。

这两种方法——Fisherian 和 Neymanian——代表了统计推断的两种哲学。理解它们的联系与区别，对于真正掌握因果推断至关重要。

\userannotation{为什么需要两种方法？}{FRT 适用于小样本精确推断和假设检验，Neymanian 适用于大样本渐近推断和置信区间构建。两者互为补充，而非相互排斥。}

\section{有限总体量}

在 Neymanian 框架中，我们需要一套精确的记号来描述\textbf{有限总体量}——这些是不带 "hat" 的真实值。

考虑一个包含 $n$ 个个体的有限总体，其中 $n_1$ 个被分配到治疗组，$n_0 = n - n_1$ 个被分配到对照组。对于单元 $i = 1, \ldots, n$，我们有潜在结果 $Y_i(1)$ 和 $Y_i(0)$，以及个体效应 $\tau_i = Y_i(1) - Y_i(0)$。

\textbf{有限总体均值}：
\[
\bar{Y}(1) = n^{-1} \sum_{i=1}^n Y_i(1), \quad
\bar{Y}(0) = n^{-1} \sum_{i=1}^n Y_i(0)
\]

\textbf{有限总体方差}（使用 $n-1$ 作为除数使定理更优雅）：
\[
S^2(1) = (n-1)^{-1} \sum_{i=1}^n \{ Y_i(1) - \bar{Y}(1) \}^2, \quad
S^2(0) = (n-1)^{-1} \sum_{i=1}^n \{ Y_i(0) - \bar{Y}(0) \}^2
\]

\textbf{有限总体协方差}：
\[
S(1,0) = (n-1)^{-1} \sum_{i=1}^n \{ Y_i(1) - \bar{Y}(1) \} \{ Y_i(0) - \bar{Y}(0) \}
\]

\textbf{个体效应}：
\[
\tau = n^{-1} \sum_{i=1}^n \tau_i = \bar{Y}(1) - \bar{Y}(0)
\]
是平均因果效应的真值。

\textbf{个体效应的方差}：
\[
S^2(\tau) = (n-1)^{-1} \sum_{i=1}^n (\tau_i - \tau)^2
\]

\userannotation{为什么需要协方差？}{因为 $Y(1)$ 和 $Y(0)$ 之间可能有相关性。如果治疗组和对照组的个体是匹配的（正相关），估计的方差会减小；如果是负相关，方差会增大。}

\section{一个关键引理}

在有限总体量的世界中，有一个优雅的关系连接了协方差与各个方差：

\begin{Lemma}[Lemma 4.1 — 方差与协方差的关系]\label{lemm:4.1}
\[
2S(1,0) = S^2(1) + S^2(0) - S^2(\tau)
\]
其中 $\tau_i = Y_i(1) - Y_i(0)$ 是个体因果效应。
\end{Lemma}

这个引理的证明是直接的代数运算。这个恒等式将在后续推导中发挥关键作用。

\textbf{推论}：如果所有个体因果效应都相等（$S^2(\tau) = 0$），那么 $S(1,0) = [S^2(1) + S^2(0)]/2$，即潜在结果完全正相关。

\section{Neyman (1923) 定理}

基于观测结果，我们可以计算样本均值：
\[
\hat{\bar{Y}}(1) = n_1^{-1} \sum_{i=1}^n Z_i Y_i, \quad
\hat{\bar{Y}}(0) = n_0^{-1} \sum_{i=1}^n (1 - Z_i) Y_i
\]

样本方差：
\[
\hat{S}^2(1) = (n_1 - 1)^{-1} \sum_{i=1}^n Z_i \{ Y_i - \hat{\bar{Y}}(1) \}^2, \quad
\hat{S}^2(0) = (n_0 - 1)^{-1} \sum_{i=1}^n (1 - Z_i) \{ Y_i - \hat{\bar{Y}}(0) \}^2
\]

\textbf{关键观察}：不存在 $S(1,0)$ 和 $S^2(\tau)$ 的样本版本，因为潜在结果 $Y_i(1)$ 和 $Y_i(0)$ 从未对任何单元同时被观测到。

现在陈述 Neyman 的核心定理：

\begin{Theorem}[Neyman (1923) 定理]\label{chap4-thm:neyman1923}
在 CRE 下：
\begin{enumerate}
  \item \textbf{无偏性}：差值均值估计量 $\hat{\tau} = \hat{\bar{Y}}(1) - \hat{\bar{Y}}(0)$ 对 $\tau$ 是无偏的：
    \[
    \mathbb{E}(\hat{\tau}) = \tau
    \]

  \item \textbf{方差公式}：$\hat{\tau}$ 有方差
    \begin{align}
    \text{var}(\hat{\tau}) &= \frac{S^2(1)}{n_1} + \frac{S^2(0)}{n_0} - \frac{S^2(\tau)}{n} \label{eq:neyman-var} \\
    &= \frac{n_0}{n_1 n} S^2(1) + \frac{n_1}{n_0 n} S^2(0) + \frac{2}{n} S(1,0) \label{eq:neyman-var-equivalent}
    \end{align}

  \item \textbf{保守方差估计}：方差估计量
    \[
    \hat{V} = \frac{\hat{S}^2(1)}{n_1} + \frac{\hat{S}^2(0)}{n_0}
    \]
    对估计 $\text{var}(\hat{\tau})$ 是保守的，即
    \[
    \mathbb{E}(\hat{V}) - \text{var}(\hat{\tau}) = \frac{S^2(\tau)}{n} \geq 0
    \]
    等号当且仅当 $\tau_i = \tau$（所有单元的个体效应恒定）时成立。
\end{enumerate}
\end{Theorem}

\textbf{关于期望和方差的含义}：潜在结果都是固定数值，只有治疗指示变量 $Z_i$ 是随机的。因此，期望和方差都相对于 $Z_i$ 的随机性——即 $n_1$ 个 1 和 $n_0$ 个 0 的所有可能排列。

\userannotation{第三项为什么是减项？}{直觉上：如果个体因果效应存在变异（$S^2(\tau) > 0$），治疗组和对照组之间的差异不仅来自随机抽样，还来自因果效应本身的真实变异。这种"真实变异"不应该被算作估计误差，所以被减去。}

\section{方差公式的理解}

方差公式 \cref{eq:neyman-var} 与经典两样本问题中的方差公式不同，因为它不仅依赖于潜在结果的有限总体方差，还依赖于个体效应的有限总体方差。

\textbf{一个令人困惑的现象}：个体效应越异质，$\hat{\tau}$ 的变异性越小。为什么？

从等价格式 \cref{eq:neyman-var-equivalent} 理解：比较正相关和负相关的潜在结果情况。

假设在一次实现的实验中，我们观察到相对较大的治疗潜在结果：
\begin{itemize}
  \item 如果 $S(1,0) > 0$，则那些接受治疗的单元也有相对较大的对照潜在结果。因此，对照下的观测结果相对较小，导致较大的 $\hat{\tau}$
  \item 如果 $S(1,0) < 0$，则那些接受治疗的单元有相对较小的对照潜在结果。因此，对照下的观测结果相对较大，导致较小的 $\hat{\tau}$
\end{itemize}

总体而言，尽管 $\hat{\tau}$ 的无偏性不依赖于潜在结果之间的相关性，但在正相关情况下更可能观察到更极端的 $\hat{\tau}$。因此，当潜在结果正相关时，$\hat{\tau}$ 的方差更大。

\section{Neyman (1923) 定理的证明}\label{sec:prove-neyman1923}

\textbf{无偏性证明}：

从表示开始：
\begin{align*}
\hat{\tau} &= n_1^{-1} \sum_{i=1}^n Z_i Y_i - n_0^{-1} \sum_{i=1}^n (1 - Z_i) Y_i \\
&= n_1^{-1} \sum_{i=1}^n Z_i Y_i(1) - n_0^{-1} \sum_{i=1}^n (1 - Z_i) Y_i(0)
\end{align*}

利用期望的线性性：
\begin{align*}
\mathbb{E}(\hat{\tau}) &= \mathbb{E}\left\{ n_1^{-1} \sum_{i=1}^n Z_i Y_i(1) - n_0^{-1} \sum_{i=1}^n (1 - Z_i) Y_i(0) \right\} \\
&= n_1^{-1} \sum_{i=1}^n \mathbb{E}(Z_i) Y_i(1) - n_0^{-1} \sum_{i=1}^n \mathbb{E}(1 - Z_i) Y_i(0) \\
&= n_1^{-1} \sum_{i=1}^n \frac{n_1}{n} Y_i(1) - n_0^{-1} \sum_{i=1}^n \frac{n_0}{n} Y_i(0) \\
&= n^{-1} \sum_{i=1}^n Y_i(1) - n^{-1} \sum_{i=1}^n Y_i(0) = \tau
\end{align*}

\textbf{方差证明}：

我们可以进一步将 $\hat{\tau}$ 写为：
\[
\hat{\tau} = \sum_{i=1}^n Z_i \left\{ \frac{Y_i(1)}{n_1} + \frac{Y_i(0)}{n_0} \right\} - n_0^{-1} \sum_{i=1}^n Y_i(0)
\]

利用简单随机抽样的引理，得到：
\begin{align*}
\text{var}(\hat{\tau}) &= \frac{n_1 n_0}{n(n-1)} \sum_{i=1}^n \left\{ \frac{Y_i(1)}{n_1} + \frac{Y_i(0)}{n_0} - \frac{\bar{Y}(1)}{n_1} - \frac{\bar{Y}(0)}{n_0} \right\}^2 \\
&= \frac{n_0}{n_1 n} S^2(1) + \frac{n_1}{n_0 n} S^2(0) + \frac{2}{n} S(1,0)
\end{align*}

利用 \cref{lemm:4.1}，我们也可以将方差写为：
\[
\text{var}(\hat{\tau}) = \frac{S^2(1)}{n_1} + \frac{S^2(0)}{n_0} - \frac{S^2(\tau)}{n}
\]

\textbf{方差估计的期望}：

因为治疗组是从 $n$ 个单元中简单随机抽取的 $n_1$ 个样本，简单随机抽样引理确保：
\[
\mathbb{E}\{\hat{S}^2(1)\} = S^2(1), \quad \mathbb{E}\{\hat{S}^2(0)\} = S^2(0)
\]

因此，$\hat{V}$ 对 \cref{eq:neyman-var} 中前两项是无偏的。

\section{渐近正态性与置信区间}

\cite{li2017} 进一步证明了基于有限总体 CLT 的 $\hat{\tau}$ 的渐近正态性：

\begin{Theorem}[CLT for 差值均值]\label{chap4-thm:clt}
令 $n \to \infty$ 和 $n_1 \to \infty$。如果 $n_1/n$ 有 $(0, 1)$ 中的极限值，$\{S^2(1), S^2(0), S(1,0)\}$ 有极限值，且
\[
\max_{1 \leq i \leq n} \{Y_i(1) - \bar{Y}(1)\}^2 / n \to 0, \quad
\max_{1 \leq i \leq n} \{Y_i(0) - \bar{Y}(0)\}^2 / n \to 0
\]
则
\[
\frac{\hat{\tau} - \tau}{\sqrt{\text{var}(\hat{\tau})}} \xrightarrow{d} N(0, 1)
\]
且
\[
\hat{S}^2(1) \to S^2(1), \quad \hat{S}^2(0) \to S^2(0) \quad \text{（依概率）}
\]
\end{Theorem}

这个定理确保当样本量足够大时，$\hat{\tau}$ 的抽样分布可以用正态分布近似。它还确保样本方差是总体方差的一致估计。

\textbf{置信区间}：这为 $\tau$ 的保守大样本置信区间提供了正当理由：
\[
\hat{\tau} \pm z_{1-\alpha/2} \sqrt{\hat{V}}
\]
其中 $z_{1-\alpha/2}$ 是标准正态分布的 $1 - \alpha/2$ 分位数。

\textbf{与弱零假设的對偶性}：通过置信区间与假设检验的对偶性，这也提供了检验
\[
H_0: \tau = 0
\]
（称为弱零假设）的基础。

\userannotation{为什么 Neymanian 方法只给出"弱零假设"的检验？}{因为在 Neymanian 框架中，我们无法观测到 $S^2(\tau)$，所以无法精确计算方差。置信区间和检验都是渐近有效的，只在样本量大时可靠。}

\section{Neymanian 推断与 OLS 的关系}

实践中，研究者经常使用基于回归的推断来估计平均因果效应。标准方法是运行 OLS 回归：
\[
(\hat{\alpha}, \hat{\beta}) = \arg\min_{(a,b)} \sum_{i=1}^n (Y_i - a - b Z_i)^2
\]
并使用 treatment 系数的估计量 $\hat{\beta}$ 作为平均因果效应的估计量。

\textbf{关键发现}：系数 $\hat{\beta}$ 恰好等于差值均值：
\[
\hat{\beta} = \hat{\tau}
\]

但是，OLS 的通常方差估计量与 $\hat{V}$ 不同。幸运的是，Eicker-Huber-White (EHW) 稳健方差估计量接近 $\hat{V}$，HC2 变体恰好等于 $\hat{V}$。

在 LaLonde 数据上的应用：
\begin{verbatim}
> tauhat = mean(y[z==1]) - mean(y[z==0])
> vhat   = var(y[z==1])/n1 + var(y[z==0])/n0
> sehat  = sqrt(vhat)
> tauhat
[1] 1794.343
> sehat
[1] 670.9967
\end{verbatim}

OLS 的标准误（通常）太小：
\begin{verbatim}
> olsfit = lm(y ~ z)
> summary(olsfit)$coef[2, 1:2]
  Estimate Std. Error
 1794.3431   632.8536
\end{verbatim}

使用 EHW 稳健标准误得到正确结果：
\begin{verbatim}
> sqrt(hccm(olsfit, type = "hc2")[2, 2])
[1] 670.9967
\end{verbatim}

\section{重尾结局与正态近似的失效}

\cref{chap4-thm:clt} 中的 CLT 在某些正则条件下成立。当潜在结局为重尾分布时，这些条件可能被违反。

当对照潜在结果被污染时（例如 Cauchy 分布），正态近似效果很差。模拟研究表明，在重尾情况下，直方图和正态近似之间存在显著差异。

\section{本章小结}

本章建立了 Neymanian 推断的基础：

\begin{enumerate}
  \item \textbf{有限总体量}：不带 hat 的真实值
    \[
    \bar{Y}(1), \bar{Y}(0), S^2(1), S^2(0), S(1,0), S^2(\tau)
    \]
  \item \textbf{Lemma 4.1}：$2S(1,0) = S^2(1) + S^2(0) - S^2(\tau)$
  \item \textbf{Neyman (1923) 定理}：
    \begin{itemize}
      \item $\mathbb{E}(\hat{\tau}) = \tau$（无偏）
      \item $\text{var}(\hat{\tau}) = \frac{S^2(1)}{n_1} + \frac{S^2(0)}{n_0} - \frac{S^2(\tau)}{n}$
      \item $\mathbb{E}(\hat{V}) = \text{var}(\hat{\tau}) + \frac{S^2(\tau)}{n}$（保守）
    \end{itemize}
  \item \textbf{方差估计}：$\hat{V} = \frac{\hat{S}^2(1)}{n_1} + \frac{\hat{S}^2(0)}{n_0}$
  \item \textbf{置信区间}：$\hat{\tau} \pm z_{1-\alpha/2} \sqrt{\hat{V}}$
  \item \textbf{CLT}：$\frac{\hat{\tau} - \tau}{\sqrt{\text{var}(\hat{\tau})}} \xrightarrow{d} N(0, 1)$
  \item \textbf{OLS 关系}：$\hat{\beta} = \hat{\tau}$，但需使用 EHW 稳健标准误
\end{enumerate}

下一章我们将学习如何通过分层（stratification）来提高估计的精度——这是一种利用协变量信息来改进推断的方法。

\section{习题}

\begin{Exercise}{\ref{ex:4.1} 引理的证明}\label{ex:4.1}
证明 Lemma \ref{lemm:4.1}：$2S(1,0) = S^2(1) + S^2(0) - S^2(\tau)$。
\end{Exercise}

\begin{Exercise}{\ref{ex:4.2} Neyman (1923) 定理的另一种证明}\label{ex:4.2}
在 CRE 下，计算：
\[
\text{var}\{\hat{\bar{Y}}(1)\}, \quad \text{var}\{\hat{\bar{Y}}(0)\}, \quad \text{cov}\{\hat{\bar{Y}}(1), \hat{\bar{Y}}(0)\}
\]
并利用这些公式计算 $\text{var}(\hat{\tau})$。
\end{Exercise}

\begin{Exercise}{\ref{ex:4.3} Neymanian 推断与 OLS}\label{ex:4.3}
证明 $\hat{\beta} = \hat{\tau}$（即 OLS 系数等于差值均值）。证明 EHW 稳健方差估计量的 HC2 变体恰好等于 $\hat{V}$。
\end{Exercise}

\begin{Exercise}{\ref{ex:4.4} 治疗效应的异质性}\label{ex:4.4}
证明 $S^2(\tau) = 0$ 蕴含 $S^2(1) = S^2(0)$。给出反例说明 $S^2(1) = S^2(0)$ 但 $S^2(\tau) \neq 0$。

证明 $S^2(1) < S^2(0)$ 蕴含 $S(Y(0), \tau) < 0$。给出反例说明 $S^2(1) > S^2(0)$ 但 $S(Y(0), \tau) < 0$。
\end{Exercise}

\begin{Exercise}{\ref{ex:4.5} 方差公式的改进上界}\label{ex:4.5}
证明方差公式的改进上界：
\[
\text{var}(\hat{\tau}) \leq \frac{1}{n}\left\{ \sqrt{\frac{n_0}{n_1}}S(1) + \sqrt{\frac{n_1}{n_0}}S(0) \right\}^2
\]
并讨论等号成立的条件。
\end{Exercise}

\begin{Exercise}{\ref{ex:4.6} Neyman (1923) 的向量版本}\label{ex:4.6}
将 Neyman (1923) 的结果推广到多结局情形。考虑向量结局 $\boldsymbol{V} \in \mathbb{R}^K$ 的平均因果效应：
\[
\tau_{\boldsymbol{V}} = \frac{1}{n}\sum_{i=1}^n \{\boldsymbol{V}_i(1) - \boldsymbol{V}_i(0)\}
\]

证明 $\hat{\tau}_{\boldsymbol{V}}$ 对 $\tau_{\boldsymbol{V}}$ 是无偏的，并求出其协方差矩阵。
\end{Exercise}

\begin{Exercise}{\ref{ex:4.7} BRE 中的推断}\label{ex:4.7}
回忆在 \cref{def:cre} 中定义的完全随机实验。在这种设定下，推导平均因果效应的方差公式，并比较其与 CRE 中方差的差异。
\end{Exercise}

\begin{Exercise}{\ref{ex:4.8} 推荐阅读}\label{ex:4.8}
\begingroup
\parskip=0pt
\textbf{推荐阅读}：\cite{ding2017} 重新审视了因果推断中的一些经典问题。\cite{imbens2015} 是关于因果推断的权威综述。

\textbf{历史注记}：Neyman 1923 年的原始论文\cite{neyman1923}奠定了有限总体推断的基础。当代统计学家将其视为因果推断范式的里程碑。
\endgroup
\end{Exercise}

% === 用户问答记录 ===
% (后续通过 interactive-qa 更新)


% Appendix
\appendix
% Appendix: Questions and Answers
% User annotations and interactive Q&A records

\appendix
\chapter{问答与注解}

\section{阅读过程中的问题与解答}

% === 用户问答记录 ===
% 本节记录阅读过程中的问答内容

\subsection*{亚组的定义}
\textbf{问}：什么是亚组？

\textbf{答}：亚组（subgroup）是由某个协变量（通常是二元变量 $X_i$）分割出来的子群体。例如按性别、年龄、收入等分割。亚组因果效应衡量的是在该亚组内的平均因果效应。

\subsection*{恒等式 $\tau = \pi_1 \tau_1 + \pi_0 \tau_0$ 的推导}
\textbf{问}：如何推导 $\tau = \pi_1 \tau_1 + \pi_0 \tau_0$？

\textbf{答}：利用指示函数的完备性 $\mathbb{I}(X_i=1) + \mathbb{I}(X_i=0) = 1$，将总体求和拆分为两个亚组求和，再除以对应的亚组样本量即可得到。

\subsection*{为什么在 $H_{0F}$ 下 $\bm{Y}$ 是固定的？}

\textbf{问}：为什么在 Fisher's sharp null ($H_{0F}: Y_i(1) = Y_i(0)$) 下，潜在结果向量 $\bm{Y}$ 是固定的？

\textbf{答}：$H_{0F}$ 断言 treatment 对每个单元都没有效应，即 $Y_i(1) = Y_i(0) = Y_i$ 对所有 $i$。这意味着整个 $\bm{Y}(1) = \bm{Y}(0) = \bm{Y}$ 向量是确定的（固定值），不再有随机性。

在 FRT 中，我们 condition on 这个固定的 $\bm{Y}$，只让治疗分配 $\bm{Z}$ 有随机性。因此，随机化分布可以精确枚举（$M = \binom{n}{n_1}$ 种可能），不依赖任何渐近假设——这就是 FRT"有限样本精确"的本质。

\textit{当你在阅读过程中提出问题时，回答将被记录在此处。}

% === 用户注解 ===
% 使用 \userannotation{问题}{解答} 格式插入

\section{学习笔记}

\textit{在这里记录你的学习心得和思考。}

\endinput


% Bibliography
\bibliographystyle{plainnat}
\bibliography{references}

\end{document}
